% --- Archon physics formalization source begin ---
% archon:physics
% archon:covers PhyXMiniProblems/problem_phyx_mini_0108.lean
% archon:source-report reports/phyx_mini/problem_phyx_mini_0108.source.json

\chapter{Physics problem phyx\_mini\_0108}
\label{ch:PhyXMiniProblems_problem_phyx_mini_0108}

\paragraph{Problem source.}
\#\# Physical scenario

A reflecting telescope (figure) is to be made by using a spherical mirror with a radius of curvature of 1.30 m and an eyepiece with a focal length of 1.10 cm. The final image is at infinity.

\#\# Question

What will the angular magnification be?

\#\# Answer choices

- A: A: 59.1

- B: B: 58.5

- C: C: 57.9

- D: D: 56.7

\#\# Figure caption (auxiliary; use the image as primary evidence)

The image is a diagram illustrating a ray tracing scenario. Here’s a detailed description:

1. **Objects and Elements**:
   - **Eye**: On the left side of the image, there is an illustration of an eye.
   - **Lens**: Positioned in the center of the diagram, represented by a thin vertical shape indicating a lens (possibly converging).
   - **Mirror**: On the right, there is a large arc suggesting the presence of a concave mirror.

2. **Rays**:
   - Several purple arrows (representing light rays) travel from the left to the right.
   - The rays converge at the lens and then extend towards the mirror.
   - After reflecting off the concave mirror, the rays are shown diverging back towards the left.

3. **Interactions**:
   - The light rays pass through the lens towards the mirror and reflect back in the direction of the eye, suggesting a focus point created possibly behind the lens.
   - The pathway indicates how light interacts through refraction at the lens and reflection at the mirror.

There is no text present in this particular diagram. The setup likely explains the basic principles of optics, such as refraction and reflection through a lens and mirror.

\#\# Dataset metadata

Optical Instruments; Multi-Formula Reasoning, Physical Model Grounding Reasoning

\paragraph{Recorded answer/context.}
A: A: 59.1

\paragraph{Figure/image path.}
/root/proposal\_for\_physic/science-mango/phyx\_mini\_run/phyx\_data/test\_image/108.png

\paragraph{Formalization target.}
create a compiling Lean file with sorry bodies at `PhyXMiniProblems/problem\_phyx\_mini\_0108.lean`. The Lean declarations must preserve the physical quantities, dimensions or dimensional roles, figure labels, governing-law hypotheses, and final relation expressed by this problem.
Use Mathlib/Physlib names found through LeanExplore where available. If a physics API is missing, introduce faithful local abstractions rather than scalar placeholder aliases.

\begin{theorem}[Physics formalization target]
\label{thm:physics:phyx_mini_0108:target}
\lean{PhyXMiniProblems.ProblemPhyXMini0108.angularMagnification_is_answer_A}
\uses{def:physics:phyx-mini-0108:phyxminiproblems-problemphyxmini0108-reflectingtelescopesetup, def:physics:phyx-mini-0108:phyxminiproblems-problemphyxmini0108-matchesproblemandfigure, def:physics:phyx-mini-0108:phyxminiproblems-problemphyxmini0108-hasphysicalparameters, def:physics:phyx-mini-0108:phyxminiproblems-problemphyxmini0108-satisfiesparaxialtelescopelaws, def:physics:phyx-mini-0108:phyxminiproblems-problemphyxmini0108-answerchoice, def:physics:phyx-mini-0108:phyxminiproblems-problemphyxmini0108-roundstonearesttenth, def:physics:phyx-mini-0108:phyxminiproblems-problemphyxmini0108-isclosestanswerchoice}
The assigned autoformalize agent should translate this physics problem into Lean declarations in the covered file, with theorem and lemma proof bodies written as `by sorry`.
\end{theorem}
\begin{proof}
This is an autoformalization task, not a proof task. Produce statements that can later be proved by the physics prover without weakening the physical model.
\end{proof}

% --- Archon declaration topology begin ---
\section{Declaration topology}
\begin{definition}[Optical Length]
  \label{def:physics:phyx-mini-0108:phyxminiproblems-problemphyxmini0108-opticallength}
  \lean{PhyXMiniProblems.ProblemPhyXMini0108.OpticalLength}
  A real-valued physical length, independent of the chosen unit readout
\end{definition}
\begin{proof}
  This is the declared model definition of Optical Length.
\end{proof}
\begin{definition}[centimeter Unit Choices]
  \label{def:physics:phyx-mini-0108:phyxminiproblems-problemphyxmini0108-centimeterunitchoices}
  \lean{PhyXMiniProblems.ProblemPhyXMini0108.centimeterUnitChoices}
  Unit choices in which length readouts are measured in centimetres
\end{definition}
\begin{proof}
  This is the declared model definition of centimeter Unit Choices.
\end{proof}
\begin{definition}[length In Meters]
  \label{def:physics:phyx-mini-0108:phyxminiproblems-problemphyxmini0108-lengthinmeters}
  \lean{PhyXMiniProblems.ProblemPhyXMini0108.lengthInMeters}
  \uses{def:physics:phyx-mini-0108:phyxminiproblems-problemphyxmini0108-opticallength}
  Read a physical length as a real number of metres
\end{definition}
\begin{proof}
  This is the declared model definition of length In Meters.
\end{proof}
\begin{definition}[length In Centimeters]
  \label{def:physics:phyx-mini-0108:phyxminiproblems-problemphyxmini0108-lengthincentimeters}
  \lean{PhyXMiniProblems.ProblemPhyXMini0108.lengthInCentimeters}
  \uses{def:physics:phyx-mini-0108:phyxminiproblems-problemphyxmini0108-opticallength, def:physics:phyx-mini-0108:phyxminiproblems-problemphyxmini0108-centimeterunitchoices}
  Read a physical length as a real number of centimetres
\end{definition}
\begin{proof}
  This is the declared model definition of length In Centimeters.
\end{proof}
\begin{definition}[Figure Landmark]
  \label{def:physics:phyx-mini-0108:phyxminiproblems-problemphyxmini0108-figurelandmark}
  \lean{PhyXMiniProblems.ProblemPhyXMini0108.FigureLandmark}
  Optical elements and focal landmark visible in the source figure
\end{definition}
\begin{proof}
  This is the declared model definition of Figure Landmark.
\end{proof}
\begin{definition}[Primary Mirror Geometry]
  \label{def:physics:phyx-mini-0108:phyxminiproblems-problemphyxmini0108-primarymirrorgeometry}
  \lean{PhyXMiniProblems.ProblemPhyXMini0108.PrimaryMirrorGeometry}
  The optical form of the primary mirror
\end{definition}
\begin{proof}
  This is the declared model definition of Primary Mirror Geometry.
\end{proof}
\begin{definition}[Eyepiece Kind]
  \label{def:physics:phyx-mini-0108:phyxminiproblems-problemphyxmini0108-eyepiecekind}
  \lean{PhyXMiniProblems.ProblemPhyXMini0108.EyepieceKind}
  The optical action of the eyepiece
\end{definition}
\begin{proof}
  This is the declared model definition of Eyepiece Kind.
\end{proof}
\begin{definition}[Telescope Ray Stage]
  \label{def:physics:phyx-mini-0108:phyxminiproblems-problemphyxmini0108-telescoperaystage}
  \lean{PhyXMiniProblems.ProblemPhyXMini0108.TelescopeRayStage}
  Successive regions of the ray path shown in the telescope diagram
\end{definition}
\begin{proof}
  This is the declared model definition of Telescope Ray Stage.
\end{proof}
\begin{definition}[Paraxial Ray Pattern]
  \label{def:physics:phyx-mini-0108:phyxminiproblems-problemphyxmini0108-paraxialraypattern}
  \lean{PhyXMiniProblems.ProblemPhyXMini0108.ParaxialRayPattern}
  Qualitative paraxial-ray patterns used to encode the diagram
\end{definition}
\begin{proof}
  This is the declared model definition of Paraxial Ray Pattern.
\end{proof}
\begin{definition}[Final Image Location]
  \label{def:physics:phyx-mini-0108:phyxminiproblems-problemphyxmini0108-finalimagelocation}
  \lean{PhyXMiniProblems.ProblemPhyXMini0108.FinalImageLocation}
  Whether the final image is finite or at optical infinity
\end{definition}
\begin{proof}
  This is the declared model definition of Final Image Location.
\end{proof}
\begin{definition}[Reflecting Telescope Setup]
  \label{def:physics:phyx-mini-0108:phyxminiproblems-problemphyxmini0108-reflectingtelescopesetup}
  \lean{PhyXMiniProblems.ProblemPhyXMini0108.ReflectingTelescopeSetup}
  \uses{def:physics:phyx-mini-0108:phyxminiproblems-problemphyxmini0108-opticallength, def:physics:phyx-mini-0108:phyxminiproblems-problemphyxmini0108-primarymirrorgeometry, def:physics:phyx-mini-0108:phyxminiproblems-problemphyxmini0108-eyepiecekind, def:physics:phyx-mini-0108:phyxminiproblems-problemphyxmini0108-telescoperaystage, def:physics:phyx-mini-0108:phyxminiproblems-problemphyxmini0108-paraxialraypattern, def:physics:phyx-mini-0108:phyxminiproblems-problemphyxmini0108-finalimagelocation}
  The physical quantities and qualitative optical data of the reflecting telescope. The requested angular magnification is stored as an unknown nonnegative magnitude, not as an answer value
\end{definition}
\begin{proof}
  This is the declared model definition of Reflecting Telescope Setup.
\end{proof}
\begin{definition}[Matches Problem And Figure]
  \label{def:physics:phyx-mini-0108:phyxminiproblems-problemphyxmini0108-matchesproblemandfigure}
  \lean{PhyXMiniProblems.ProblemPhyXMini0108.MatchesProblemAndFigure}
  \uses{def:physics:phyx-mini-0108:phyxminiproblems-problemphyxmini0108-lengthinmeters, def:physics:phyx-mini-0108:phyxminiproblems-problemphyxmini0108-lengthincentimeters, def:physics:phyx-mini-0108:phyxminiproblems-problemphyxmini0108-reflectingtelescopesetup}
  Problem-text and primary-figure readouts. The mirror is concave and spherical; the eyepiece is converging; incident light and light delivered to the eye are parallel, with reflection through the primary focus between them
\end{definition}
\begin{proof}
  This is the declared model definition of Matches Problem And Figure.
\end{proof}
\begin{definition}[Has Physical Parameters]
  \label{def:physics:phyx-mini-0108:phyxminiproblems-problemphyxmini0108-hasphysicalparameters}
  \lean{PhyXMiniProblems.ProblemPhyXMini0108.HasPhysicalParameters}
  \uses{def:physics:phyx-mini-0108:phyxminiproblems-problemphyxmini0108-lengthinmeters, def:physics:phyx-mini-0108:phyxminiproblems-problemphyxmini0108-reflectingtelescopesetup}
  Positivity conditions for the telescope's physical parameters
\end{definition}
\begin{proof}
  This is the declared model definition of Has Physical Parameters.
\end{proof}
\begin{definition}[Satisfies Paraxial Telescope Laws]
  \label{def:physics:phyx-mini-0108:phyxminiproblems-problemphyxmini0108-satisfiesparaxialtelescopelaws}
  \lean{PhyXMiniProblems.ProblemPhyXMini0108.SatisfiesParaxialTelescopeLaws}
  \uses{def:physics:phyx-mini-0108:phyxminiproblems-problemphyxmini0108-reflectingtelescopesetup}
  Governing paraxial-optics laws. For a spherical mirror, the radius of curvature is twice its focal length. At normal adjustment (final image at infinity), the magnitude of a telescope's angular magnification is the primary-to-eyepiece focal-length ratio. These relations contain no numerical magnification or answer-choice conclusion
\end{definition}
\begin{proof}
  This is the declared model definition of Satisfies Paraxial Telescope Laws.
\end{proof}
\begin{definition}[Answer Choice]
  \label{def:physics:phyx-mini-0108:phyxminiproblems-problemphyxmini0108-answerchoice}
  \lean{PhyXMiniProblems.ProblemPhyXMini0108.AnswerChoice}
  Labels of the displayed multiple-choice answers
\end{definition}
\begin{proof}
  This is the declared model definition of Answer Choice.
\end{proof}
\begin{definition}[angular Magnification Readout]
  \label{def:physics:phyx-mini-0108:phyxminiproblems-problemphyxmini0108-answerchoice-angularmagnificationreadout}
  \lean{PhyXMiniProblems.ProblemPhyXMini0108.AnswerChoice.angularMagnificationReadout}
  \uses{def:physics:phyx-mini-0108:phyxminiproblems-problemphyxmini0108-answerchoice}
  Dimensionless angular-magnification magnitude printed beside each choice
\end{definition}
\begin{proof}
  This is the declared model definition of angular Magnification Readout.
\end{proof}
\begin{definition}[Rounds To Nearest Tenth]
  \label{def:physics:phyx-mini-0108:phyxminiproblems-problemphyxmini0108-roundstonearesttenth}
  \lean{PhyXMiniProblems.ProblemPhyXMini0108.RoundsToNearestTenth}
  Agreement with a value displayed to the nearest tenth
\end{definition}
\begin{proof}
  This is the declared model definition of Rounds To Nearest Tenth.
\end{proof}
\begin{definition}[Is Closest Answer Choice]
  \label{def:physics:phyx-mini-0108:phyxminiproblems-problemphyxmini0108-isclosestanswerchoice}
  \lean{PhyXMiniProblems.ProblemPhyXMini0108.IsClosestAnswerChoice}
  \uses{def:physics:phyx-mini-0108:phyxminiproblems-problemphyxmini0108-answerchoice}
  A displayed choice is strictly closest to the physical magnification
\end{definition}
\begin{proof}
  This is the declared model definition of Is Closest Answer Choice.
\end{proof}
\begin{lemma}[primary Mirror Focal Length In Meters]
  \label{lem:physics:phyx-mini-0108:phyxminiproblems-problemphyxmini0108-primarymirrorfocallengthinmeters}
  \lean{PhyXMiniProblems.ProblemPhyXMini0108.primaryMirrorFocalLengthInMeters}
  \uses{def:physics:phyx-mini-0108:phyxminiproblems-problemphyxmini0108-lengthinmeters, def:physics:phyx-mini-0108:phyxminiproblems-problemphyxmini0108-reflectingtelescopesetup, def:physics:phyx-mini-0108:phyxminiproblems-problemphyxmini0108-matchesproblemandfigure, def:physics:phyx-mini-0108:phyxminiproblems-problemphyxmini0108-hasphysicalparameters, def:physics:phyx-mini-0108:phyxminiproblems-problemphyxmini0108-satisfiesparaxialtelescopelaws}
  The 1.30 m spherical mirror has paraxial focal length 0.65 m
\end{lemma}
\begin{proof}
  The conclusion follows from the governing relations and figure readouts named in the statement of primary Mirror Focal Length In Meters.
\end{proof}
% --- Archon declaration topology end ---
% --- Archon physics formalization source end ---
